\documentclass[11pt]{beamer}
%\usepackage{CJKutf8}
\usepackage{beamerthemesplit}
\usepackage{}
%\usetheme{Frankfurt}
%\usetheme{CambridgeUS}
%\usecolortheme{beaver}
%\usetheme{AnnArbor}
\usetheme{Dresden}
%\usebeamercolor{beetle}
\usepackage{xeCJK}
\setCJKmainfont{AR PL KaitiM GB}
%\useoutertheme{miniframes}
\usepackage{amsmath}
\usepackage{graphicx}
\usepackage{float} 
\usepackage{subfigure}
\usepackage{amssymb}
\usepackage{graphicx}
\usepackage{eufrak}
\usepackage{color}
\usepackage{array}
\usepackage{slashed}
\usepackage{simplewick}
\usepackage{tikz}
\usepackage{tcolorbox}
\usepackage[T1]{fontenc}
\graphicspath{{../figures/}}

%%figures
\def\lfig#1#2{\includegraphics[width=#1 in]{#2}}
\def\addfig#1#2{\begin{center}\includegraphics[width=#1 in]{#2}\end{center}}
\def\wulian{\includegraphics[width=0.18in]{emoji_wulian.jpg}}
\def\bigwulian{\includegraphics[width=0.35in]{emoji_wulian.jpg}}
\def\bye{\includegraphics[width=0.18in]{emoji_bye.jpg}}
\def\bigbye{\includegraphics[width=0.35in]{emoji_bye.jpg}}
\def\huaixiao{\includegraphics[width=0.18in]{emoji_huaixiao.jpg}}
\def\bighuaixiao{\includegraphics[width=0.35in]{emoji_huaixiao.jpg}}
\def\jianxiao{\includegraphics[width=0.18in]{emoji_jianxiao.jpg}}
\def\bigjianxiao{\includegraphics[width=0.35in]{emoji_jianxiao.jpg}}
%% colors
\def\blacktext#1{{\color{black}#1}}
\def\bluetext#1{{\color{blue}#1}}
\def\redtext#1{{\color{red}#1}}
\def\darkbluetext#1{{\color[rgb]{0,0.2,0.6}#1}}
\def\skybluetext#1{{\color[rgb]{0.2,0.7,1.}#1}}
\def\cyantext#1{{\color[rgb]{0.,0.5,0.5}#1}}
\def\greentext#1{{\color[rgb]{0,0.7,0.1}#1}}
\def\darkgray{\color[rgb]{0.2,0.2,0.2}}
\def\lightgray{\color[rgb]{0.6,0.6,0.6}}
\def\gray{\color[rgb]{0.4,0.4,0.4}}
\def\blue{\color{blue}}
\def\red{\color{red}}
\def\green{\color{green}}
\def\darkgreen{\color[rgb]{0,0.4,0.1}}
\def\darkblue{\color[rgb]{0,0.2,0.6}}
\def\skyblue{\color[rgb]{0.2,0.7,1.}}
%%control
\def\be{\begin{equation}}
\def\ee{\nonumber\end{equation}}
\def\bea{\begin{eqnarray*}}
\def\eea{\nonumber\end{eqnarray*}}
\def\bch{}
\def\ech{}
\def\bitem{\begin{itemize}}
\def\eitem{\end{itemize}}
\def\bcenter{\begin{center}}
\def\ecenter{\end{center}}
\def\bex{\begin{minipage}{0.2\textwidth}\includegraphics[width=0.6in]{jugelizi.png}\end{minipage}\begin{minipage}{0.76\textwidth}}
\def\eex{\end{minipage}}
\def\chtitle#1{\frametitle{\bch#1\ech}}
\def\bmat#1{\left(\begin{array}{#1}}
\def\emat{\end{array}\right)}
\def\bcase#1{\left\{\begin{array}{#1}}
\def\ecase{\end{array}\right.}
\def\bmini#1{\begin{minipage}{#1\textwidth}}
\def\emini{\end{minipage}}
\def\tbox#1{\begin{tcolorbox}#1\end{tcolorbox}}
\def\pfrac#1#2#3{\left(\frac{\partial #1}{\partial #2}\right)_{#3}}
%%symbols
\def\bropt{\,(\ \ \ )}
\def\sone{$\star$}
\def\stwo{$\star\star$}
\def\sthree{$\star\star\star$}
\def\sfour{$\star\star\star\star$}
\def\sfive{$\star\star\star\star\star$}
\def\rint{{\int_\leftrightarrow}}
\def\roint{{\oint_\leftrightarrow}}
\def\stdHf{{\textit{\r H}_f}}
\def\deltaH{{\Delta \textit{\r H}}}
\def\ii{{\dot{\imath}}}
\def\skipline{{\vskip0.1in}}
\def\skiplines{{\vskip0.2in}}
\def\lagr{{\mathcal{L}}}
\def\hamil{{\mathcal{H}}}
\def\vecv{{\mathbf{v}}}
\def\vecx{{\mathbf{x}}}
\def\vecy{{\mathbf{y}}}
\def\veck{{\mathbf{k}}}
\def\vecp{{\mathbf{p}}}
\def\vecn{{\mathbf{n}}}
\def\vecA{{\mathbf{A}}}
\def\vecP{{\mathbf{P}}}
\def\vecsigma{{\mathbf{\sigma}}}
\def\hatJn{{\hat{J_\vecn}}}
\def\hatJx{{\hat{J_x}}}
\def\hatJy{{\hat{J_y}}}
\def\hatJz{{\hat{J_z}}}
\def\hatj#1{\hat{J_{#1}}}
\def\hatphi{{\hat{\phi}}}
\def\hatq{{\hat{q}}}
\def\hatpi{{\hat{\pi}}}
\def\vel{\upsilon}
\def\Dint{{\mathcal{D}}}
\def\adag{{\hat{a}^\dagger}}
\def\bdag{{\hat{b}^\dagger}}
\def\cdag{{\hat{c}^\dagger}}
\def\ddag{{\hat{d}^\dagger}}
\def\hata{{\hat{a}}}
\def\hatb{{\hat{b}}}
\def\hatc{{\hat{c}}}
\def\hatd{{\hat{d}}}
\def\hatN{{\hat{N}}}
\def\hatH{{\hat{H}}}
\def\hatp{{\hat{p}}}
\def\Fup{{F^{\mu\nu}}}
\def\Fdown{{F_{\mu\nu}}}
\def\newl{\nonumber \\}
\def\vece{\mathrm{e}}
\def\calM{{\mathcal{M}}}
\def\calT{{\mathcal{T}}}
\def\calR{{\mathcal{R}}}
\def\barpsi{\bar{\psi}}
\def\baru{\bar{u}}
\def\barv{\bar{\upsilon}}
\def\qeq{\stackrel{?}{=}}
\def\torder#1{\mathcal{T}\left(#1\right)}
\def\rorder#1{\mathcal{R}\left(#1\right)}
\def\contr#1#2{\contraction{}{#1}{}{#2}#1#2}
\def\trof#1{\mathrm{Tr}\left(#1\right)}
\def\trace{\mathrm{Tr}}
\def\comm#1{\ \ \ \left(\mathrm{used}\ #1\right)}
\def\tcomm#1{\ \ \ (\text{#1})}
\def\slp{\slashed{p}}
\def\slk{\slashed{k}}
\def\calp{{\mathfrak{p}}}
\def\veccalp{\mathbf{\mathfrak{p}}}
\def\Tthree{T_{\tiny \textcircled{3}}}
\def\pthree{p_{\tiny \textcircled{3}}}
\def\dbar{{\,\mathchar'26\mkern-12mu d}}
\def\erf{\mathrm{erf}}
\def\const{\mathrm{constant}}
\def\pheat{\pfrac p{\ln T}V}
\def\vheat{\pfrac V{\ln T}p}
%%units
\def\fdeg{{^\circ \mathrm{F}}}
\def\cdeg{^\circ \mathrm{C}}
\def\atm{\,\mathrm{atm}}
\def\angstrom{\,\text{\AA}}
\def\SIL{\,\mathrm{L}}
\def\SIkm{\,\mathrm{km}}
\def\SIyr{\,\mathrm{yr}}
\def\SIGyr{\,\mathrm{Gyr}}
\def\SIV{\,\mathrm{V}}
\def\SImV{\,\mathrm{mV}}
\def\SIeV{\,\mathrm{eV}}
\def\SIkeV{\,\mathrm{keV}}
\def\SIMeV{\,\mathrm{MeV}}
\def\SIGeV{\,\mathrm{GeV}}
\def\SIcal{\,\mathrm{cal}}
\def\SIkcal{\,\mathrm{kcal}}
\def\SImol{\,\mathrm{mol}}
\def\SIN{\,\mathrm{N}}
\def\SIHz{\,\mathrm{Hz}}
\def\SIm{\,\mathrm{m}}
\def\SIcm{\,\mathrm{cm}}
\def\SIfm{\,\mathrm{fm}}
\def\SImm{\,\mathrm{mm}}
\def\SInm{\,\mathrm{nm}}
\def\SImum{\,\mathrm{\mu m}}
\def\SIJ{\,\mathrm{J}}
\def\SIW{\,\mathrm{W}}
\def\SIkJ{\,\mathrm{kJ}}
\def\SIs{\,\mathrm{s}}
\def\SIkg{\,\mathrm{kg}}
\def\SIg{\,\mathrm{g}}
\def\SIK{\,\mathrm{K}}
\def\SImmHg{\,\mathrm{mmHg}}
\def\SIPa{\,\mathrm{Pa}}
\def\secpage#1#2{\begin{frame}\bch\bcenter{\bf \Huge #1} \skipline \tbox{\bcenter #2\ecenter}\ecenter\ech\end{frame}}

\graphicspath{{figures/}}


\newcommand{\field}{\mathscr{F}}

\newcommand{\reals}{\mathbb{R}}
\newcommand{\complexs}{\mathbb{C}}
\newcommand{\ints}{\mathbb{Z}}
%\newcommand{\dim}{\mathrm{dim\ }}
\newcommand{\diag}{\mathrm{diag \ }}
\newcommand{\up}{\uparrow}
\newcommand{\down}{\downarrow}
\newcommand{\su}{\mathfrak{su}}
\newcommand{\so}{\mathfrak{so}}
\newcommand{\tr}{\mathrm{tr\ }}
\newcommand{\card}{\mathrm{card \ }}
\newcommand{\lag}{\mathcal{L}}
\newcommand{\hamiltonian}{\mathcal{H}}
\newcommand{\op}{\mathcal{O}}
\newtheorem{thm}{定理}[section]
\newtheorem{axm}{公理}[section]
\newtheorem{dfn}{定义}[section]
\newtheorem{experience}{经验}[section]


%\cpic{<尺寸>}{<文件名>}}用于生成居中的图片。
\newcommand{\cpic}[2]{
\begin{center}
\includegraphics[scale=#1]{#2}
\end{center}
}

%\cpicn{<尺寸>}{<文件名>}{<注释>}用于生成居中且带有注释的图片，其label为图片名。
\newcommand{\cpicn}[3]
{
\begin{figure}[h!]
\cpic{#1}{#2}
\caption{#3\label{#2}}
\end{figure}
}

\title{Quantum Field Theory\\ 正则量子化}
  \author{Haoting Xu}
  \date{\today}


\begin{document}

\begin{frame}
 
\maketitle

%\vskip 0.2in
\begin{center}
xuht9@mail2.sysu.edu.cn
\vskip 0.1in
{\tiny \url{https://github.com/HaotingXu/seminar_lec} }\\
\end{center}
\end{frame}
\section{单粒子量子力学的终结}
\secpage{单粒子量子力学的终结}{$P\simeq e^{-m|x|}$}

\begin{frame}\frametitle{海森堡绘景}
薛定谔绘景：算符不随时间改变，态随时间改变。
\be
|\psi (t)\rangle_S = e^{-i\hat{H}t}|\psi(0)\rangle_S
\ee 
海森堡绘景：算符随时间改变，态不随时间改变。为了保持算符平均值不变，定义
\be
\hat{\op}_H = e^{-i\hat{H}t} \hat{\op}_S e^{i\hat{H}t}
\ee
利用$i\partial_t |\psi(t)\rangle_S = \hat{H}|\psi(0)\rangle_S$，可以得到算符的演化方程
\be
\frac{d\hat{\op}_H(t)}{dt} = -i\left[\hat{\op}_H(t),\hat{H}\right] 
\ee
\end{frame}
\begin{frame}\frametitle{海森堡绘景的好处}
如果采用了海森堡绘景，算符不随时间变化。即算符是洛伦兹不变的。在我们讨论相对论的时候，这就good
\end{frame}
\begin{frame}\frametitle{单粒子量子力学的终结}
我们现在要试图把单粒子量子力学和狭义相对论结合起来，即简单的设想$E_{\vec{p}} = \sqrt{\vec{p}^2+m^2}$，我们采取归一化$\langle \vec{x}|\vec{p}\rangle$。狭义相对论要求无法和具有类空距离的事件沟通，即要求两个态不能相互跃迁。因此我们来计算(其中$t>|\vec{x}|=x$)
\be
\mathcal{A} = \langle \vec{x}| e^{-i\hat{H}t}|\vec{x}=0\rangle
\ee
利用套路，有
\bea
\langle \vec{x}| e^{-i\hat{H}t}|\vec{x}=0\rangle &=& \int d^3 p \langle \vec{x}|e^{-i\hat{H}t}|\vec{p}\rangle\langle\vec{p}|\vec{x}=0\rangle\\
&=& \int d^3 p \frac{1}{(2\pi)^3}e^{i\vec{p}\cdot\vec{x}}e^{-iE_pt}
\eea
\end{frame}
\begin{frame}{进行积分}
一阵优秀的化直角坐标为球坐标的操作后，得到
\be
\mathcal{A} = \frac{-i}{(2\pi)^2 x} \int_{-\infty}^{\infty} dp p e^{ipx} e^{-it\sqrt{p^2+m^2}}
\ee
进行围到积分最终得到
\be
\mathcal{A} \simeq e^{-m|\vec{x}|}
\ee
因此我们可以看到，仍然是有概率跃迁到类空距离的点，这样会违背因果律。所以单粒子的量子力学体系想要和相对论统一起来是不可能的。
\end{frame}
\begin{frame}{量子场论对于因果律的要求}
我们构建的量子场论不能破坏因果律，在量子场论中，我们把场看做算符，这意味着要求
\be
\left[\hat{\op}(x),\hat{\op}(y)\right] = 0 
\ee
如果$(x-y)^2<0$。
\end{frame}
\section{对称性与诺特定理}
\secpage{对称性与诺特定理}{一种对称性意味着守恒流守恒}
\begin{frame}\frametitle{对称性}
\begin{dfn}[对称性]
如果在一种无穷小变换$\delta \phi$下，拉氏密度$\lag$差某个函数的四维散度
\be
\delta \lag = \partial_\mu W^\mu
\ee
则称这样的变换是一种对称变换。
\end{dfn}
如果这样，我们有
\be
\delta \lag = \frac{\partial \lag}{\partial \phi}\delta \phi +  \frac{\partial \lag}{\partial (\partial_\mu \phi)} \delta (\partial_\mu \phi)
\ee
如果场还满足欧拉-拉格朗日方程，则有
\be
\partial_\mu \left(\frac{\partial \lag}{\partial(\partial_\mu \phi)}\delta \phi\right) = \partial_\mu W^\mu
\ee
\end{frame}
\begin{frame}\frametitle{诺特定理}
于是我们有
\be
\partial_\mu j^\mu =0 
\ee
其中守恒流为
\be
j^\mu = \frac{\partial \lag}{\partial(\partial_\mu \phi)}\delta \phi - W^\mu
\ee
我们便证明了诺特定理。
\end{frame}
\begin{frame}\frametitle{一个例子：能量动量张量}
考虑时空平移的无穷小变换
\be
x^\nu \rightarrow x^\nu - \epsilon^\nu
\ee
则拉氏密度作如下变换
\be
\lag(x)\rightarrow \lag(x) + \epsilon^\nu \partial_\nu \lag(x)
\ee
代入诺特定理，得到时空平移的守恒流为
\be
j^\mu = \frac{\partial \lag}{\partial(\partial_\mu \phi)}\epsilon^\nu\partial_\nu \phi  -\epsilon^\mu \lag(x)
\ee
我们可以去掉$\epsilon^\nu$，从而提取出一个张量
\be
(j^\mu)_\nu =\frac{\partial \lag}{\partial(\partial_\mu \phi)}\partial_\nu \phi  -\delta^\mu_\nu \lag(x)\equiv T^\mu_\nu
\ee
这就是能量动量张量。对应于在经典力学里我们老生常谈的:时间平移不变性对应能量守恒，空间平移不变性对应动量守恒。
\end{frame}
\begin{frame}\frametitle{小练习:KG场的能量动量张量}
套用上面的公式，求出KG场的能量动量张量
\be
\lag = \frac{1}{2} \eta^{\mu\nu} \partial_\mu \phi \partial_\nu \phi -\frac{1}{2}m^2\phi^2
\ee
\end{frame}
\begin{frame}\frametitle{角动量守恒}
大家回去自己尝试推导在无穷小洛伦兹变换下
\be
\Lambda^T \eta \Lambda = \eta
\ee
的守恒流。
\end{frame}
\begin{frame}\frametitle{内禀对称性}
例如对于复数场
\be
\lag = \partial_\mu \psi^{*}\partial^\mu \psi - V(|\psi|^2)
\ee
拉氏密度在如下变换下保持不变
\be
\psi \rightarrow \psi e^{i\alpha}
\ee
这种对称性被叫做$U(1)$对称性，试着得到$U(1)$对称性的守恒流。一会我们将看到这个守恒流的意义。
\end{frame}
\section{正则量子化}
\begin{frame}\frametitle{正则量子化一般步骤}
我们现在要建立场的量子化描述，一般步骤如下
\begin{enumerate}	
	\item  写出拉氏密度。
	\item  写出哈密顿密度。$\hamiltonian = \pi^a\cdot \dot{\phi}_a-\lag$
	\item 引入对易关系。
	\item 将场使用产生湮灭算符展开，代入哈密顿量。
	\item  Normal Ordering.
\end{enumerate}
\end{frame}
\begin{frame}\frametitle{一个例子：标量场量子化}
我们先来搞个最简单的
\be
\lag = \frac12 \partial^\mu\phi\partial_\mu \phi-\frac12 m^2 \phi^2
\ee
瞬间得到哈密顿量
\be
\hamiltonian = \frac12 \dot{\phi}^2 + \frac12 (\nabla \phi)^2+ \frac12 m^2\phi^2
\ee
将场变为算符，引入对易关系
\be
\left[\hat{\phi}(t,\vec{x}),\hat{\pi}^0(t,\vec{y})\right] = i\delta^{(3)}(\vec{x}-\vec{y})
\ee
\end{frame}
\begin{frame}\frametitle{相对论归一化}
我们希望我们接下来书写的一些表达式都是相对论不变的。显然，下面的式子
\be
\int d^4p \delta(0) = \int d^4p \delta(p_0^2 - \vec{p}^2-m^2) 
\ee
是洛伦兹不变的，利用厂主神奇公式$\delta(f(x)) = \sum_n \frac{\delta(x-x_n)}{f^{\prime}(x_n)}$，得到
\be
\int d^4p \delta(p_0^2 - \vec{p}^2-m^2) = \int \frac{d^3p}{2E_{\vec{p}}}
\ee
是洛伦兹不变的。
\end{frame}
\begin{frame}\frametitle{相对论归一化}
我们还可以得到$2E_{\vec{p}} \delta^{3} (\vec{p}-\vec{q})$是洛伦兹不变的。我们还可以定义相对论归一化的态$|p\rangle = \sqrt{2E_{\vec{p}}}|\vec{p}\rangle$，他们满足
\be
\langle p|q\rangle = 2E_{\vec{p}} \delta^{(3)}(\vec{p}-\vec{q})
\ee
\end{frame}
\begin{frame}\frametitle{产生湮灭算符}
可以将场用产生湮灭算符展开
\be
\hat{\phi}(\vec{x}) = \int \frac{d^3p}{(2\pi)^{\frac32}}\frac{1}{\sqrt{2E_{\vec{p}}}}(\hat{a}_{\vec{p}}e^{i\vec{p}\cdot\vec{x}}+\hat{a}^\dagger_{\vec{p}}e^{-i\vec{p}\cdot\vec{x}})
\ee
如果我们考虑海森堡绘景
\be
\hat{\phi}(x) = \hat{U}^\dagger \hat{\phi}(\vec{x})\hat{U}=e^{i\hat{H}t} \hat{\phi}(\vec{x})e^{-i\hat{H}t}
\ee
先计算$\hat{U}^\dagger a_{\vec{p}} \hat{U}$
\end{frame}
\begin{frame}
利用对易关系$[H,\hat{a}_{\vec{p}}]=-E_{\vec{p}}a_{\vec{p}}$
和公式
\be
e^{\hat{A}}\hat{B}e^{-\hat{A}} = \hat{B}+ \left[\hat{A},\hat{B}\right]+ \frac{1}{2!}\left[ A,\left[\hat{A},\hat{B}\right]\right]+\cdots
\ee
我们得到
\be
e^{i\hat{H}t}a_{\vec{p}}e^{-i\hat{H}t} = a_{\vec{p}}e^{-iE_{\vec{p}} t}
\ee
\end{frame}
\begin{frame}\frametitle{最终形式}
最终我们得到
\be
\hat{\phi}(\vec{x}) = \int \frac{d^3p}{(2\pi)^{\frac32}}\frac{1}{\sqrt{2E_{\vec{p}}}}(\hat{a}_{\vec{p}}e^{-ip\cdot x}+\hat{a}^\dagger_{\vec{p}}e^{ip\cdot x})
\ee
另外从对易关系我们可以得到升降算符的对易
\end{frame}
\begin{frame}\frametitle{哈密顿量}
将上面的表达式带入到哈密顿量中
\be
H = \frac{1}{2} \int d^3x \dot{\phi}^2 + (\nabla \phi)^2+m^2 \phi^2
\ee
进行一顿计算后得到
\be
H = \int d^3p E_{\vec{p}} \left(\hat{a}^\dagger_{\vec{p}}\hat{a}_{\vec{p}} + \frac{1}{2}\delta^{(3)} (0)\right)
\ee
可见最后一项必然导致发散，那么怎么处理呢？
\end{frame}
\begin{frame}\frametitle{Normal Ordering}
有一种解释就是我们只关心能量差，因此直接把后面那项扔掉。我们只关心产生算符在前，湮灭算符在后的项，这种产生在前，湮灭在后的顺序就叫做Normal Ordering，用$N$来表示Normal Ordering，有
\be
N[\hat{A}\hat{B}\hat{C}^\dagger\cdots \hat{Z}] = (\text{creation operators on the left})\equiv :\hat{A}\hat{B}\hat{C}^\dagger\cdots \hat{Z}:
\ee
对于玻色场，上面的ordering没有问题。注意对于费米场，会额外多一个排序因子$(-1)^P$。
\end{frame}
\begin{frame}\frametitle{Casimir Effect}
你以为这样扔掉就完了吗？没有粒子的真空那部分能量不会体现吗？考虑一个长为$L$的盒子，考虑中间有两块平行板$d$。将哈密顿量拿过来
\be
H = \int d^3p  E_{\vec{p}} \frac{1}{2}\delta^{(3)} (0)
\ee
回忆$\delta$函数相当于空间的体积
\be
(2\pi)^3 \delta^{(3)}(0) = \lim_{L\rightarrow \infty}\int_{-L/2}^{L/2}d^3x e^{i\vec{p}\cdot\vec{x}}\arrowvert_{\vec{p}=0}=V
\ee
因此蕴含在两板之间的能量
\be
E\sim \frac{n\pi}{2d}
\ee
\end{frame}
\begin{frame}\frametitle{Ultraviolet cut-off}
但是两板之间的能量仍然是无穷的，实验中一般使用电磁场，我们知道电磁场在高频的时候就不很好的遵循低频边界条件了，所以我们直接让高频的部分能量衰减。
\be
E(d)\sim \sum_n \frac{n\pi}{2d} e^{-n\pi a/d}\simeq \frac{d}{2\pi a^2} - \frac{\pi}{24d}+\cdots
\ee
真正的能量是
\be
E = E(d)+ E(L-d)
\ee
最终能计算出受力
\be
F = -\frac{\partial E}{\partial d} = -\frac{\pi}{24d^2}
\ee
这就是Casimir效应，这告诉我们两个板放在一起会损失一些模式，从而表现为负能量。另一个解释是这个力起源于真空中的涨落(正反粒子对的产生和湮灭)。
\end{frame}

\begin{frame}\frametitle{Feynman's Interpreatation of the negative frequency}
\be
\hat{\phi}(\vec{x}) = \int \frac{d^3p}{(2\pi)^{\frac32}}\frac{1}{\sqrt{2E_{\vec{p}}}}(\hat{a}_{\vec{p}}e^{-ip\cdot x}+\hat{a}^\dagger_{\vec{p}}e^{ip\cdot x})
\ee
盯着这个表达式，你会发现第一项似乎是“负能量态”，如何给这一个解释？费曼给出了一个猜想，负能量态就是产生一个反粒子向外跑，而正能量态是正粒子向内跑湮灭。对于这个标量场，它是自身的反粒子。
\end{frame}
\section{Complex Field}
\secpage{复数场}{$\lag = \partial^\mu \psi^\dagger(x)\partial_\mu \psi(x)-m^2\psi^\dagger(x)\psi(x)$}
\begin{frame}\frametitle{复数场的量子化}
复数场的拉氏密度
\be
\lag = \partial^\mu \psi^\dagger(x)\partial_\mu \psi(x)-m^2\psi^\dagger(x)\psi(x)
\ee
得到哈密顿密度
\be
\hamiltonian =\partial_0 \psi^\dagger (x)\partial_0 \psi(x)+ \nabla \psi^\dagger(x)\cdot \nabla \psi(x) + m^2 \psi^\dagger(x)\psi(x)
\ee
引入对易关系
\be
\left[\hat{\psi}(t,\vec{x}),\hat{\pi}_\psi(t,\vec{y})\right]=\left[\hat{\psi}^\dagger(t,\vec{x}),\hat{\pi}_{\psi^\dagger}(t,\vec{y})\right]=i\delta^{(3)}(\vec{x}-\vec{y})\ee
\end{frame}
\begin{frame}{正反粒子}
将场算符用产生湮灭算符展开
\be
\hat{\psi}(x) = \int \frac{d^3p}{(2\pi)^{3/2}}\frac{1}{\sqrt{2E_{\vec{p}}}} \left( \hat{a}_{\vec{p}}e^{-ip\cdot x} +  \hat{b}^\dagger_{\vec{p}} e^{ip\cdot x}\right)
\ee
其中$\hat{a}_{\vec{p}},\hat{b}^\dagger_{\vec{p}}$分别湮灭两种不同的粒子，又由于这两种粒子具有相同的能量，所以我们把它们解释为正反粒子。
\end{frame}
\begin{frame}{哈密顿量}
将上述展开带入哈密顿量中，并进行normal ordering，得到
\be
\hat{H} = \int d^3p E_p \left(a^\dagger_p a_p+b^\dagger_pb_p\right)
\ee
\end{frame}
\begin{frame}{守恒流}
回顾之前$U(1)$对称性导致的复数场的守恒流
\be
J^\mu = i[(\partial^\mu \psi^\dagger)\psi-(\partial^\mu\psi)\psi^\dagger]
\ee
考虑守恒荷$j^0$，将之前的展开带入，进行normal ordering 得到
\be
N[\hat{Q}] = \int d^3 p (b^\dagger_pb_p-a^\dagger_pa_p)
\ee
使用粒子数算符，有
\be
Q = -N_a+N_b
\ee
这说明粒子数减去反粒子数守恒。这就是电荷守恒。一般来说，对正粒子选取正号，对于反粒子选取负号，即上面的整个式子差个负号。
\end{frame}
\begin{frame}\frametitle{非相对论极限}
非相对论情形下能量为
\be
E = mc^2 + \epsilon
\ee
其中$\epsilon$是一个小量，取非相对论极限的步骤就是将$mc^2$的演化部分分离出来，即
\be
\phi(x)\rightarrow \psi (x)e^{-imt}
\ee
例如对于KG方程$(\partial^2+m^2)\psi(x)e^{-imt}$，我们将上面的式子代入并做一些近似得到
\be
i \hbar\partial_t \psi = -\frac{\hbar^2}{2m}\nabla^2 \psi
\ee
即自由粒子的薛定谔方程。
\end{frame}
\begin{frame}\frametitle{复数场的情形}
对于复数场，我们采取如下分离方式
\be
\psi = \frac{1}{\sqrt{2m}} e^{-imt} \Psi
\ee
将这个带入到拉氏量中，得到
\be
\lag = i \Psi^\dagger(x)\partial_0(x) -\frac{1}{2m}\nabla \Psi^\dagger(x)\cdot \Psi(x) -\frac{g}{2}[\Psi^\dagger(x)\Psi(x)]^2
\ee
试着将这个场量子化，如果懒，就看看书上Example 12.4。书上还给出了另一种有趣的量子化方式，大家可以随便看看。
\end{frame}
\section{很多分量的场}
\secpage{很多分量的场}{升降算符数目取决于场的独立分量数}
\begin{frame}{同位旋}
1932年，海森堡发现质子的质量和中子的质量十分相近$m_p\simeq m_n$，并认为他们是同一个东西，它们组成一个类似于自旋$1/2$的系统，在一种特定的变换下可以相互转换。我们知道$SU(2)\simeq SO(3)$，我们先来考虑一个$SO(3)$的同位旋系统。假设$SO(3)$的weights分别对应与三种粒子$(t,d,h)$(short for Tom, Dick, Harry)。因为场是粒子的激发态，所以我们猜测需要三个场来描述。
\end{frame}
\begin{frame}\frametitle{拉氏量}
于是我们的场为$\vec{\Phi}(x) = (\phi_1(x),\phi_2(x),\phi_3(x))$，随便构造一个$SO(3)$(对场变换)的拉氏量
\be
\lag = \frac12 (\partial^\mu \vec{\Phi})\cdot (\partial_\mu \vec{\Phi}) - \frac{m^2}{2}\vec{\Phi}\cdot \vec{\Phi}
\ee
现在将场量子化，首先计算哈密顿密度
\be
\hamiltonian = \sum_\alpha \left[\frac12 (\partial_0 \phi_\alpha)^2+\frac12 (\nabla \phi_\alpha)^2+\frac12 m^2 \phi_\alpha^2\right]
\ee
\end{frame}
\begin{frame}\frametitle{对易关系和展开}
因此可以引入对易关系，由于我们认为三个场是独立的，有
\be
[\Phi_\alpha(x),\pi_\beta(y)] = i \delta^3 (\vec{x}-\vec{y})\delta_{\alpha\beta}
\ee
所以需要引入三组产生湮灭算符，将场展开为
\be
\vec{\Phi}(x) = \int \frac{d^3p}{(2\pi)^{3/2}} \sum_{\alpha=1}^3 \vec{h}_\alpha \left(a_{\vec{p}\alpha}e^{-ip\cdot x}+a^\dagger_{\vec{p}\alpha}e^{ip\cdot x}  \right)
\ee
其中$\vec{h}_1 = (1,0,0),\cdots$
\end{frame}
\begin{frame}\frametitle{诺特定理}
回忆三维空间的无穷小转动
\be
\begin{pmatrix}
1 &\theta_3 &-\theta_2 \\
-\theta^3 &1 &\theta_1 \\
\theta_2 & -\theta_1 & 1
\end{pmatrix}
\begin{pmatrix}\phi_1 \\\phi_2 \\ \phi_3 \end{pmatrix}
\ee
可以写为
\be
\Phi_i \rightarrow \Phi_i - \epsilon_{ijk}\theta_j \Phi_k
\ee
使用诺特定理，得到守恒流
\be
(j^\mu)_j = \epsilon_{jik} (\partial^\mu \phi_i)\phi_k
\ee
这里的守恒荷就是同位旋。
\end{frame}
\section{Guage Field}
\begin{frame}\frametitle{规范对称性}
考虑复数场
\be
\lag = (\partial^\mu \psi)^\dagger(\partial_\mu \psi)- m^2 \psi^\dagger\psi
\ee
作$U(1)$变换$\psi\rightarrow \psi e^{i\alpha}$，拉氏量不变。但是如果$\alpha$依赖于时空坐标$\alpha(x)$，那么不变性就不成立。这就说明这个场论不是全局对称的，而是局域对称的。为了使得不变性成立，引入一个新的场并定义一个新的协变导数算符
\be
D_\mu = \partial_\mu +iq A_\mu(x)
\ee
其中$A_\mu$按照如下规则变换
\be
A_\mu \rightarrow A_\mu - \frac{1}{q}\partial_\mu \alpha(x)
\ee
可以证明，在变换下$D(\psi)\rightarrow D(\psi e^{i\alpha})$，所以拉氏量改写为
\be
\lag = (D^\mu\psi)^\dagger (D_\mu \psi) -m^2\psi^\dagger\psi
\ee
\end{frame}
\begin{frame}\frametitle{规范场的拉氏量}
对于规范场，引入拉氏量
\be
\lag =-\frac{1}{4}F_{\mu\nu}F^{\mu\nu}-A^\mu J^\mu
\ee
运动方程就是麦克斯韦方程组。如果进行规范变换
\be
A_\mu(x) \rightarrow A_\mu(x) - \partial_\mu \chi(x)
\ee
则不会改变观测量。可见场的选择有一定任意性，可以进行规范固定，如选取洛伦兹规范
\be
\partial_\mu A^\mu =0
\ee
但是这样还不能完全固定场，还可以加上一个$\partial^2\epsilon=0$的场。如果再选取库伦规范$\nabla \cdot \vec{A}=0$，有了这两个规范，场的独立分量就是$4-1-1=2$个，意味着光子有两种偏振态。
\end{frame}
\begin{frame}\frametitle{场量子化}
拉氏量
\be
\lag = -\frac{1}{4} F_{\mu\nu}F^{\mu\nu}
\ee
求得哈密顿量
\be
\hamiltonian = \frac{1}{2}(\vec{E}^2+ \vec{B}^2)
\ee
引入对易关系,为了保证规范不变，我们取
\be
[A^i(x),E^j(y)] = i \int \frac{d^3p}{(2\pi)^3}e^{i(\vec{x}-\vec{y})}\left(\delta^{ij}- \frac{p^ip^j}{p^2}\right) = i \delta_{\rm tr}^{(3)}(\vec{x}-\vec{y})
\ee
仍然有$[\hat{a}_{\vec{p}\lambda},\hat{a}^\dagger_{\vec{p}\lambda}]=\delta^{(3)}(\vec{p}-\vec{q})\delta_{\lambda\lambda^{\prime}}$
\end{frame}

\begin{frame}
将场展开为
\be
A^\mu (x) = \int \frac{d^3p}{(2\pi)^{3/2}} \sum_{\lambda=1}^2 \left(\epsilon_\lambda^\mu a_{\vec{p}\lambda}e^{-ip\cdot x}+\epsilon_\lambda^{\mu *} a^\dagger_{\vec{p}\lambda}e^{ip\cdot x} \right)
\ee
将这个代入哈密顿量，据说就能得到
\be
\hat{H}= \int d^3p \sum_\lambda E_{\vec{p}}a^\dagger_{\vec{p}\lambda} a_{\vec{p}\lambda}
\ee
解释为光子。
\end{frame} 
\end{document}
